% !TEX program = lualatex
% 11: 情報モラルと社会：技術の進歩と責任・リスク

\documentclass[aspectratio=43,professionalfonts,handout,10pt]{beamer}
\usepackage{beamer_template}

\newcommand{\LectureNum}{11}
\newcommand{\LectureTitle}{情報モラルと社会：技術の進歩と責任・リスク}
\newcommand{\CourseName}{情報リテラシー}
\newcommand{\TextPages}{10-11}

\begin{document}
% --- Slide 1: title ---

\begin{frame}{\CourseName\quad 第\LectureNum 回「\LectureTitle 」}
  {\small \TermName\quad \TeacherName\quad ／\quad 教科書 \TextPages}

  \vfill

  \begin{block}{今日の流れ}
    \small
    \begin{columns}[T]
      \begin{column}{0.48\textwidth}
        \begin{enumerate}
          \item 電子商取引（EC）
          \item 人工知能（AI）の影響
          \item VR/AR・ナビ・IoT
        \end{enumerate}
      \end{column}
      \begin{column}{0.48\textwidth}
        \begin{enumerate}\setcounter{enumi}{3}
          \item 健康への影響
          \item デジタルデバイド
          \item 生成AI
        \end{enumerate}
      \end{column}
    \end{columns}
  \end{block}
  \vfill

  \begin{block}{今日の目標}
    \begin{itemize}
      \item 情報技術の発展に伴う社会的責任やリスクについて説明できる。
    \end{itemize}
  \end{block}
  \vfill

  \begin{exampleblock}{キーワード}
    \small \textbf{EC} ／ \textbf{電子マネー} ／ \textbf{深層学習} ／ \textbf{VR/AR} ／ \textbf{IoT} ／ \textbf{VDT症候群} ／ \textbf{デジタルデバイド} ／ \textbf{生成AI}
  \end{exampleblock}
\end{frame}

% --- Slide 2: 電子商取引 ---

\begin{frame}{電子商取引（教科書 p.22）}
  \begin{block}{ECの分類}
    \small
    \begin{itemize}
      \item \textbf{B to B}（Business to Business）：企業間のオンライン取引
      \item \textbf{B to C}（Business to Consumer）：企業から消費者（Amazon・楽天等）
      \item \textbf{C to C}（Consumer to Consumer）：個人間売買（メルカリ・ヤフオク等）
      \item \textbf{電子マネー}：デジタル情報で通貨の機能を果たす（Suica・PayPay等）
      \item \textbf{ネットオークション}：C to Cの形式。入札で売買
    \end{itemize}
  \end{block}

  \begin{exampleblock}{演習}
    自分が使ったことのあるEC（B to CまたはC to C）を1つ挙げ，便利な点・不安な点を話し合ってみよう。
  \end{exampleblock}

  \begin{slidenote}
    「誰から誰に」で分類を覚えよう。電子マネーは支払い手段，ECは取引の形態という違いも押さえておこう。
  \end{slidenote}
\end{frame}

% --- Slide 3: 人工知能の影響 ---

\begin{frame}{人工知能（AI）の影響（教科書 p.22）}
  \begin{block}{AIとは}
    \small
    コンピュータに人間のような知的な動作をさせる技術。\\
    \textbf{深層学習（ディープラーニング）}により，大量のデータから高精度な認識・予測が可能になった。
  \end{block}

  \begin{exampleblock}{主な応用分野}
    \small
    \begin{itemize}
      \item 顔認証による本人確認
      \item 医療における診断支援
      \item 災害予測・自動運転
      \item ルーティン業務の自動化→雇用への影響
    \end{itemize}
  \end{exampleblock}

  \begin{slidenote}
    深層学習は人間の脳神経回路を模した技術。身近な例として，スマホの顔認証や音声アシスタント（Siri・Alexa等）がある。AIの進化は雇用構造も変えつつある。
  \end{slidenote}
\end{frame}

% --- Slide 4: VR/AR・ナビ・IoT ---

\begin{frame}{VR/AR・ナビゲーション・IoT（教科書 p.23）}
  \begin{columns}[T]
    \begin{column}{0.32\textwidth}
      \begin{block}{VR / AR / MR}
        \small
        \begin{itemize}
          \item \textbf{VR}：完全な仮想空間
          \item \textbf{AR}：現実に仮想情報を重ねる
          \item \textbf{MR}：現実と仮想が相互作用
        \end{itemize}
      \end{block}
    \end{column}
    \begin{column}{0.32\textwidth}
      \begin{exampleblock}{ナビゲーション}
        \small
        \begin{itemize}
          \item GPS衛星の電波を受信
          \item 遅延時間から現在位置を計算
          \item カーナビ・地図アプリ等
        \end{itemize}
      \end{exampleblock}
    \end{column}
    \begin{column}{0.32\textwidth}
      \begin{exampleblock}{IoT}
        \small
        \begin{itemize}
          \item 身の回りのモノがネットに接続
          \item センサーを活用したシステムが普及
          \item スマート家電・交通・医療等
        \end{itemize}
      \end{exampleblock}
    \end{column}
  \end{columns}

  \begin{slidenote}
    VR＝仮想のみ，AR＝現実＋仮想，MR＝両者が相互作用。ARの例：ポケモンGO。GPSは4機以上の衛星で位置を特定。IoTとAIを組み合わせたスマートホーム・スマートシティも広がっている。
  \end{slidenote}
\end{frame}

% --- Slide 6: 健康への影響 ---

\begin{frame}{健康への影響（教科書 p.24）}
  \begin{block}{VDT症候群}
    \small
    ディスプレイ（VDT）を長時間使用することで生じる健康障害（眼精疲労・目の疲れ等）。正しい姿勢と適度な休憩で軽減できる。
  \end{block}

  \begin{columns}[T]
    \begin{column}{0.48\textwidth}
      \begin{exampleblock}{ネット依存}
        \small
        スマートフォンがないと不安になり，常に操作してしまう状態
      \end{exampleblock}
    \end{column}
    \begin{column}{0.48\textwidth}
      \begin{exampleblock}{テクノストレス}
        \small
        IT機器の操作に対して生じるストレス（不安・依存の2種類）
      \end{exampleblock}
    \end{column}
  \end{columns}

  \begin{exampleblock}{演習}
    VDT症候群を防ぐために，自分が日頃できる対策を話し合ってみよう。
  \end{exampleblock}

  \begin{slidenote}
    VDT＝Visual Display Terminal。20〜20〜20ルール（20分ごとに20フィート先を20秒見る）が推奨されている。テクノストレスはIT機器が苦手な人（不安症）と使いすぎる人（依存症）どちらにも現れる。
  \end{slidenote}
\end{frame}

% --- Slide 7: デジタルデバイド ---

\begin{frame}{デジタルデバイドへの対応（教科書 p.24）}
  \begin{block}{デジタルデバイドとは}
    所得・年齢・地域などの違いによって，得られる情報の質や量，発信できる機会に格差が生じること。
  \end{block}

  \vfill

  \begin{alertblock}{問題点}
    情報格差は\textbf{社会的・経済的格差をさらに広げる}。\\
    生活に必要な情報を公平・安全に提供できる環境の整備が必要。
  \end{alertblock}

  \begin{exampleblock}{演習}
    デジタルデバイドの具体例を1つ挙げ，解消のために何ができるか話し合ってみよう。
  \end{exampleblock}

  \begin{slidenote}
    例：高齢者のスマホ操作困難・地方の通信インフラ不足・発展途上国のネット接続率の低さ。公共Wi-Fi整備・ICT教育・バリアフリー端末の普及などが対策として挙げられる。
  \end{slidenote}
\end{frame}

% --- Slide 8: 生成AI ---

\begin{frame}{生成AI（教科書 p.25）}
  \begin{block}{生成AIとは}
    \small
    学習したデータをもとに，新しいコンテンツ（\textbf{文章・画像・動画・音楽}等）を自動的に生み出す人工知能。\\
    ユーザーの指示（プロンプト）に応じてコンテンツを生成する。
  \end{block}

  \vfill

  \begin{alertblock}{使用上の注意}
    \small
    \begin{itemize}
      \item 生成された情報が正しいとは限らない（ハルシネーション）
      \item 著作権・個人情報の取り扱いに注意
      \item 入力した情報が学習データとして利用される場合がある
    \end{itemize}
  \end{alertblock}

  \begin{slidenote}
    代表例：ChatGPT（文章）・Midjourney（画像）・Suno（音楽）等。生成AIは「ツール」であり，使う人の判断力・リテラシーが重要になる。
  \end{slidenote}
\end{frame}

% --- チェックリスト ---

\begin{frame}{まとめ・チェックリスト}
  \begin{block}{自己チェック（できたら $\checkmark$ をつけよう）}
    \begin{itemize}
      \item[$\square$] 電子商取引（B to B/C/C）と電子マネーの違いを説明できる
      \item[$\square$] AIの応用例（顔認証・医療・自動運転等）を挙げられる
      \item[$\square$] VR・AR・MRの違いを説明できる
      \item[$\square$] VDT症候群・ネット依存・テクノストレスの違いを理解した
      \item[$\square$] デジタルデバイドの意味と情報格差の問題を理解した
      \item[$\square$] 生成AIの仕組みと使用上の注意点を理解した
    \end{itemize}
  \end{block}

  \vfill

  {\small 質問があれば授業後またはTeamsで受け付けます。}
\end{frame}

\end{document}
